\documentclass[12pt]{article}
\usepackage[T1]{fontenc}
\usepackage[utf8]{inputenc}
\usepackage{lmodern}
\usepackage{textcomp}
\usepackage{enumitem}
\usepackage{geometry}
\geometry{margin=1in}
\begin{document}
\begin{center}
Answer Key - Mathematics - K-12 Level Assessment 13 \
\small (Answers are ordered exactly as in the question paper.)
\end{center}

\section*{Multiple Choice}
\begin{enumerate}[label=\arabic*.]
\item \textbf{Answer:} B
\\ \textit{Options:} A) -2; B) -4; C) -1; D) 1
\\ \textit{Note:} The polynomial has degree 2 only if the cubic and quadratic terms from the expansion cancel each other out, which occurs when \( p = -4 \) and \( q = 0 \), leading to \( p + q = -4 \).
\item \textbf{Answer:} B
\\ \textit{Options:} A) 75; B) 76; C) 22; D) 23
\\ \textit{Note:} The correct answer is 76 because the list includes all integers from 25 to 100, and to find the total count, we can subtract the smallest number from the largest and then add one, resulting in 100 - 25 + 1 = 76.
\item \textbf{Answer:} C
\\ \textit{Options:} A) 2500; B) 2; C) 50; D) 25
\\ \textit{Note:} The area of the rectangle is 50 because if the width is $w$, then the length is $2 w$, and using the Pythagorean theorem with the diagonal $5\sqrt{5}$ leads to the equation $(2 w)^2 + w^2 = (5\sqrt{5})^2$, which simplifies to $5 w^2 = 125$, giving $w^2 = 25$, and thus the area is $2 w \cdot w = 50$.
\item \textbf{Answer:} C
\\ \textit{Options:} A) 11; B) 13; C) 8; D) 5
\\ \textit{Note:} The correct answer is 8 because by applying Fermat's Little Theorem, which states that if p is a prime and a is an integer not divisible by p, then \( a^{p-1} \equiv 1 \mod p \), we find that \( 25 \equiv 2 \mod 23 \) and thus \( 2^{22} \equiv 1 \mod 23 \), allowing us to calculate \( 25^{1059} \mod 23 \) as \( 2^{1059 \mod 22} = 2^{17} \equiv 8 \mod 23 \).
\item \textbf{Answer:} C
\\ \textit{Options:} A) 41; B) 22140; C) 19270; D) 2870
\\ \textit{Note:} The correct answer is 19270 because it utilizes the formula for the sum of squares, specifically calculating the sum from 1 to 40 and subtracting the sum from 1 to 20, which yields the total for the squares from 21 to 40.
\end{enumerate}

\section*{True / False}
\begin{enumerate}[label=\arabic*.]
\setcounter{enumi}{5}
\item \textbf{Answer:} False
\\ \textit{Options:} A) True; B) False
\\ \textit{Note:} The total number of liters of water in 8 fish tanks, each with 40 liters, is actually 320 liters, not 280 liters, since 8 multiplied by 40 equals 320.
\item \textbf{Answer:} False
\\ \textit{Options:} A) True; B) False
\\ \textit{Note:} The solution to the equation -2 k = -34.8 is k = 17.4, not k = -17.4, because dividing both sides of the equation by -2 gives a positive result.
\item \textbf{Answer:} False
\\ \textit{Options:} A) True; B) False
\\ \textit{Note:} The area of the square is 36 square units, as the radius of the circle is 6 units (derived from the circumference of 12π), resulting in a square area of 6² = 36.
\item \textbf{Answer:} False
\\ \textit{Options:} A) True; B) False
\\ \textit{Note:} The statement is false because 4 over 9 of 72 students is actually 32 students, not 49.
\item \textbf{Answer:} True
\\ \textit{Options:} A) True; B) False
\\ \textit{Note:} Sapphire can make 25 bouquets using 7 balloons each because 25 multiplied by 7 equals 175, which is less than her total of 179 balloons.
\end{enumerate}

\section*{Long Form Response}
\begin{enumerate}[label=\arabic*.]
\setcounter{enumi}{10}
\item \textbf{Answer:} To calculate the minimum typing speed needed for Marshall to complete his English paper on time, we need to determine how many words he must type per minute. He has 550 words to type and 25 minutes available. By dividing the total number of words (550) by the total time in minutes (25), we find that he needs to type at least 22 words per minute. This means if Marshall types at this speed or faster, he will complete his paper before class starts.
\\ \textit{Note:} correct because it accurately applies the formula for calculating typing speed by dividing the total number of words by the total time available, resulting in a minimum requirement of 22 words per minute to ensure completion on time.
\item \textbf{Answer:} In mathematics, an angle in relation to a circle is a measure of how much the circle is turned around its center point. A full circle consists of 360 degrees, which means that each degree represents 1/360 of the entire circle. To determine the measure of an angle that represents 1/360 of a full circle, we simply recognize that this angle is equal to 1 degree. This division of the circle into 360 equal parts allows us to describe angles in a standardized way, making it easier to work with geometric concepts. Thus, an angle measuring 1/360 of a full circle is precisely 1 degree.
\\ \textit{Note:} correct because it accurately defines the relationship between angles and circles by stating that a full circle is divided into 360 degrees, thereby establishing that an angle representing 1/360 of a circle is equal to 1 degree.
\end{enumerate}

\vfill
\noindent\textit{End of Answer Key}
\end{document}